% Ad Atticum 3.7 (ad-atticum-03-07) --- English translation
% Translated by Claude (Anthropic) from the Latin (Perseus).
% Style guide: STYLE.md.

\ciceroLetterOpener{CICERO ATTICO salutem}

\ciceroSection{1} I came to Brundisium on the fourteenth day before the Kalends of May. There your slaves brought me a letter from you; and other slaves on the third day after that brought another letter. As to your asking and urging that I should be with you in Epirus, your goodwill is most pleasing to me and not at all new. The plan would indeed be the one I should wish, if I were allowed to spend all my time there; for I hate crowds, I shun men, I can scarcely look upon the daylight; that solitude would be, especially in so familiar a place, not bitter to me. But to turn aside thither for the sake of my journey: first, it is out of the way; next, four days' march from Autronius and the rest; and lastly, without you. For a fortified place would profit a man dwelling there, but is not necessary for one passing through. If I dared, I should make for Athens; the chance was so falling out that I should have wished it. Now both our enemies are there, and we have not you, and we fear that they may interpret even that town as not sufficiently far from Italy; nor do you write to what day we are to expect you.

\ciceroSection{2} Your calling me back to life accomplishes the one thing, that I keep my hands from myself; the other you cannot accomplish, that I should not be sorry for my decision and my life. For what is there to hold me, especially if the hope which attended me as I was setting out is no more? I shall not enumerate all the miseries I have fallen into through the highest injury and crime --- not so much of my enemies as of those who envied me --- lest I both stir up my own grief and call you into the same mourning. This I affirm: that no one has ever been afflicted with so great a calamity, nor has anyone had death so much to be longed for. The most honourable moment for meeting it has been let go; the times that remain are not for medicine but for the end of pain.

\ciceroSection{3} On the commonwealth, I see you gathering everything that you think can bring me some hope of a change of affairs. Slight as these are, since you wish it, let us wait. Do you nonetheless overtake us if you have hastened; for either we shall come up to Epirus, or we shall go slowly through Candavia. The doubt about Epirus came not from any fickleness of mine, but because we did not know where we should see my brother; whom indeed I do not know how I shall see, nor how I shall let go again. That is the greatest and most wretched of all my miseries. I should write to you both more often and at greater length, but my grief has taken from me, of all the parts of my mind, this faculty most. I long to see you. Take care of your health. Sent the day before the Kalends of May, from Brundisium.
